\documentclass[11pt,letterpaper]{article}
\usepackage[margin=1in]{geometry}
\usepackage[T1]{fontenc}
\usepackage{lmodern,amsmath,amssymb,amsthm,mathtools,booktabs}
\usepackage[colorlinks=true,linkcolor=blue,citecolor=blue,urlcolor=blue]{hyperref}
\usepackage{microtype}
\newcommand{\R}{\mathbb R}
\newcommand{\E}{\mathbb E}
\newcommand{\KL}{\operatorname{KL}}
\newcommand{\Cov}{\operatorname{Cov}}
\newcommand{\softmax}{\operatorname{softmax}}
\newcommand{\tr}{\operatorname{tr}}
\newtheorem{proposition}{Proposition}[section]
\title{Separate Comparison and Predictive Content\\Whole-Belief Q/K Readouts and General Output Mixing}
\author{Working construction for the gauge--VFE program}
\date{September 14, 2026}
\begin{document}
\maketitle
\begin{abstract}
Comparison and prediction need not update the same random variable. This distinction permits independent query and key measurements to determine source weights while separately learned content and output maps determine what is returned. We develop this construction with whole-belief inputs, retain the full covariance of every reduced Gaussian comparison, and distinguish a predictive mixture mean from a precision-weighted belief update. A normalized conditional model recovers the usual four-map attention head under an explicit key-norm prior, whereas uncertain-comparison KL routing yields a restricted variational construction. General output mixing permits each attention head to contribute to every output coordinate without requiring dense cross-block gauge transport. We give the transformation laws, parameter counts, gradient paths, implementation selectors, and limits of the current one-pass implementation. The architectural correspondence does not establish equivalence of training objectives, Gaussian belief dynamics, or language-model performance.
\end{abstract}
\tableofcontents

\section{The question and its three distinct changes}
The proposal combines three decisions that should be assessed separately. First, comparison selects sources and predictive content supplies the returned quantity. Second, each attention head may read the entire canonical belief through its own independent query, key, and content maps. Third, the output of every head may be mixed into all coordinates of the returned representation. Removing a restriction on input access does not by itself change a belief-update law, and introducing independent values does not by itself give a general output map.

The source construction already appears in the live working manuscript \path{relational_gauge_attention.tex}, particularly its sections on separate conditional comparison and content, uncertain comparison, and whole-belief input scope. The same distinction between an exact predictive construction and a same-query gradient flow appears in \texttt{correspondence.tex}; the separate-frame and learned-value discussions in \path{GL(K)_attention.tex} supply related geometric context. This paper develops those constructions together and records the corresponding implementation. These manuscripts provide research continuity, not independent certification of a claim. Every identity used below has its assumptions stated where it is derived.

The motivation is architectural. In standard multi-head attention, learned per-head projections act on the whole input, weighted value vectors are concatenated, and a general learned output projection acts on that concatenation \cite{vaswani}. The present changes reproduce those input and output access patterns. They leave the existing Gaussian KL router in place. No experiment here establishes that an observed perplexity plateau was caused by the older restrictions or that the new path improves it.

\section{Beliefs, frames, and attention-head coordinates}
Let token $i$ carry a local Gaussian belief with mean $\mu_i\in\R^K$ and covariance $\Sigma_i\succ0$. Let $A_i,B_i\in GL(K)$ be realized query and key frames. Tied frames mean $A_i=B_i=U_i$. Define canonical role moments by
\begin{equation}
 \xi_i^Q=A_i^{-1}\mu_i,\quad C_i^Q=A_i^{-1}\Sigma_i A_i^{-\top},\qquad
 \xi_i^K=B_i^{-1}\mu_i,\quad C_i^K=B_i^{-1}\Sigma_i B_i^{-\top}.
 \label{eq:canonical}
\end{equation}
A frame is a change of coordinates; a rectangular measurement below is generally information losing. These operations should not be identified.

The geometric partition consists of blocks indexed by $g$, with $K=\sum_g d_g$. A compact block-diagonal frame stores the $d_g\times d_g$ blocks of $A_i$ and $B_i$ without allocating a dense $K\times K$ matrix. Attention heads are indexed by $h$. They designate separate comparison and content channels, rather than necessarily designating disjoint subsets of the random belief. The general mathematics permits different numbers of gauge blocks and attention heads. The current implementation uses $H$ equal gauge blocks and $H$ attention heads with $d=K/H$.

For each head choose independent matrices
\begin{equation}
 S_Q^h,S_K^h\in\R^{r_h\times K},\qquad
 T_h\in\R^{s_h\times K},\qquad O_h\in\R^{K\times s_h}.
 \label{eq:maps}
\end{equation}
Query and key dimensions must agree within a comparison head. The content dimension need not equal the comparison dimension. The learned content and output maps may be rank deficient, because the predictive mean requires only their linear action. A nonsingular Gaussian KL in $r_h$ dimensions, however, requires that both comparison covariances be positive definite. Full row rank of $S_Q^h,S_K^h$, together with positive definite $C_i^Q,C_j^K$, suffices.

\section{Whole-belief inputs with independent Q/K and learned T/O}
\subsection{Exact reduced Gaussian comparisons}
The role-specific Gaussian marginals are
\begin{align}
 \widehat q_i^h&=\mathcal N(a_i^h,G_i^{Q,h}), &
 a_i^h&=S_Q^h\xi_i^Q, &G_i^{Q,h}&=S_Q^h C_i^Q(S_Q^h)^\top,\\
 \widehat k_j^h&=\mathcal N(b_j^h,G_j^{K,h}), &
 b_j^h&=S_K^h\xi_j^K, &G_j^{K,h}&=S_K^h C_j^K(S_K^h)^\top.
 \label{eq:marginals}
\end{align}
Each reduced covariance is a full $r_h\times r_h$ matrix even when the local $\Sigma_i$ was diagonal. Replacing it by its diagonal changes the comparison distribution. The independent matrices need not be orthogonal, share a row space, or define an invertible transport between the original beliefs.

For $\delta_{ij}^h=b_j^h-a_i^h$, the directed Gaussian divergence is
\begin{equation}
 D_{ij}^h=\frac12\left[
 \tr\big((G_j^{K,h})^{-1}G_i^{Q,h}\big)
 +(\delta_{ij}^h)^\top(G_j^{K,h})^{-1}\delta_{ij}^h-r_h
 +\log\frac{\det G_j^{K,h}}{\det G_i^{Q,h}}\right].
 \label{eq:kl}
\end{equation}
With a nonempty allowed source set $\mathcal J_i$, positive temperature $\tau_h$, and finite source log-bias $\ell_{ij}^h$ on allowed entries, routing is
\begin{equation}
 \beta_{ij}^h=
 \frac{\exp(\ell_{ij}^h-D_{ij}^h/\tau_h)}
 {\sum_{k\in\mathcal J_i}\exp(\ell_{ik}^h-D_{ik}^h/\tau_h)},\qquad j\in\mathcal J_i.
 \label{eq:router}
\end{equation}
Forbidden entries receive zero weight. A row with no admitted source is outside this normalized-mixture derivation and must follow an explicit software convention. The implementation also retains its configured upper KL clamp; a saturated numerical router uses the clamped energy rather than exactly \eqref{eq:kl}. All variational identities stated for the Gaussian divergence assume that this clamp is inactive.

Independent full-rank square Q/K maps remain useful learned operators: invertibility does not make their relative action disappear. KL invariance applies when the same invertible change of comparison coordinates acts on both distributions. Two different learned maps do not meet that condition. Conversely, learning arbitrary matrices introduces conditioning and identifiability issues that an orthonormal shared measurement avoids.

\subsection{What whole-belief access changes}
Let $E_h\in\R^{K\times d}$ insert a vector into gauge block $h$. A head-local map has the form
\begin{equation}
 S_Q^h=\overline S_Q^h E_h^\top,\qquad
 S_K^h=\overline S_K^h E_h^\top,\qquad
 T_h=\overline T_h E_h^\top.
 \label{eq:local}
\end{equation}
Whole-belief input scope removes this column-support restriction. A change to any input block may then alter the query, key, or content of any head. It does not require changing the gauge group: multiplying a rectangular map by a compact block frame can be evaluated block by block. It also does not increase the total dimensionality of the belief. A $K$-coordinate token already had $K$ coordinates before the restriction was removed.

The head-local family is contained in the full-input family by zeroing off-block columns. The inclusion concerns realizable functions, not equal initialization, optimization speed, or equal parameter count. In the implementation, full-input Q/K maps begin with independent random orthonormal rows across the full input width. Content maps begin with the previous local identity or projector embedded in the wider matrix, with off-block columns zero but trainable. Thus the initialization is conservative for content and different for comparison.

\subsection{Overlapping readouts are not independent random variables}
Several heads may measure the same Gaussian draw. For tied canonical covariance $C_i$,
\begin{equation}
 \Cov(S_Q^h X_i,S_Q^l X_i)=S_Q^h C_i(S_Q^l)^\top.
 \label{eq:crosshead}
\end{equation}
This generally does not vanish. Each marginal Gaussian comparison remains valid, but summing marginal KL divergences does not produce the joint KL of the stacked readouts unless the necessary factorization assumptions hold. If the stacked measurement is redundant, its joint covariance can even be singular. Independent routing channels are an architectural choice; they do not prove independence of measurements of one underlying belief.

\section{Separate comparison and predictive content}
\subsection{A normalized finite-source model}
Fix a conditioning context consisting of vectors $q_i,k_j\in\R^r$, content vectors $v_j\in\R^s$, biases $b_{ij}$, temperature $\tau>0$, and $R\succ0$. The context is fixed while computing the conditional distribution; its parameters may still receive gradients during outer training. On a finite nonempty source set define
\begin{equation}
 Z_i^{\rm prior}=\sum_j\exp\left(b_{ij}+\frac{\|k_j\|^2}{2\tau}\right),\qquad
 \pi_{ij}=\frac{\exp(b_{ij}+\|k_j\|^2/(2\tau))}{Z_i^{\rm prior}}.
\end{equation}
Let $J$ denote the selected source, $Y$ a comparison observation, and $Z$ predictive content. The joint law
\begin{equation}
 p_i(j,y,z)=\pi_{ij}\mathcal N(y;k_j,\tau I_r)\mathcal N(z;v_j,R)
 \label{eq:joint}
\end{equation}
normalizes by integrating the two conditional Gaussian densities and summing $\pi_{ij}$. It stipulates $Y\perp Z\mid J$ at fixed context. The means may be calculated from a common deterministic representation without violating that conditional independence.

Expanding the Gaussian comparison exponent at $Y=q_i$ gives
\begin{equation}
 b_{ij}+\frac{\|k_j\|^2}{2\tau}
 -\frac{\|q_i-k_j\|^2}{2\tau}
 =b_{ij}+\frac{q_i^\top k_j}{\tau}-\frac{\|q_i\|^2}{2\tau}.
\end{equation}
The last term is constant across sources. Therefore
\begin{equation}
 p_i(J=j\mid Y=q_i)=\softmax_j\left(\frac{q_i^\top k_j}{\tau}+b_{ij}\right),\qquad
 \E[Z\mid Y=q_i]=\sum_j p_i(j\mid q_i)v_j.
 \label{eq:exactattention}
\end{equation}
Choosing independent linear maps for $q_i,k_j,v_j$ and applying an output map reproduces the four-map attention function. The compensating key-norm prior is essential to this exact Gaussian construction; replacing it by a uniform source prior leaves a negative key-norm term. This identity does not establish that the architecture was originally derived from this model or that its usual training objective optimizes the comparison evidence.

Indeed, integrating out content yields
\begin{equation}
 -\log p_i(q_i)=\frac r2\log(2\pi\tau)+\frac{\|q_i\|^2}{2\tau}
 +\log Z_i^{\rm prior}-\log\sum_j e^{b_{ij}+q_i^\top k_j/\tau}.
\end{equation}
At fixed context this contains neither the values nor the output map. Learning useful T/O maps therefore requires a content-sensitive objective, such as downstream prediction. Conditioning source inference on an observed downstream target would generally alter the posterior weights; an ordinary forward attention pass is not automatically that target-conditioned posterior.

\subsection{Uncertain comparison and KL routing}
For a supplied receiver distribution $\widehat q_i(y)$, normalized key densities $p_j^S(y)$, normalized content densities $r_j(z)$, and a positive prior $\pi_j$, consider
\begin{equation}
 p(j,y,z)=\pi_jp_j^S(y)r_j(z),\qquad
 r_\beta(j,y,z)=\beta_j\widehat q_i(y)r_j(z).
\end{equation}
Canceling the content terms inside the log ratio and integrating gives
\begin{equation}
 \KL(r_\beta\|p)=\sum_j\beta_j\KL(\widehat q_i\|p_j^S)+\KL(\beta\|\pi).
 \label{eq:restricted}
\end{equation}
At temperature one, minimizing over the simplex gives \eqref{eq:router} with $\ell_j=\log\pi_j$. More generally, the strictly convex categorical objective
\begin{equation}
 \mathcal F_\tau(\beta)=\sum_j\beta_jD_j+\tau\KL(\beta\|\pi)
\end{equation}
has stationary equation $D_j+\tau\log(\beta_j/\pi_j)+\tau+\lambda=0$, and hence the same temperature-scaled routing law. Arbitrary temperature is an entropy-weighted variational objective; it is not the unmodified KL in \eqref{eq:restricted}. The optimum value is $-\tau\log\sum_j\pi_j e^{-D_j/\tau}$.

The content mean under this restricted family is $\sum_j\beta_j\E_{r_j}Z$. This provides a normalized uncertain-comparison construction for the existing projected-KL router. It is a restricted variational optimum with a supplied receiver distribution, whereas \eqref{eq:exactattention} is an exact conditional posterior for a point comparison observation. Neither statement identifies arbitrary updates of the receiver belief with optimization of a shared global free energy.

\subsection{Why a shared latent draw changes content}
A different model uses one random canonical belief $X\mid J=j\sim\mathcal N(\xi_j,C_j)$ and observes
\begin{equation}
 Y=SX+\varepsilon,\qquad V=TX,\qquad
 \varepsilon\sim\mathcal N(0,R_y),\qquad R_j^Y=SC_jS^\top+R_y\succ0.
\end{equation}
The cross covariance is $TC_jS^\top$. Gaussian conditioning then gives
\begin{align}
 \E[V\mid Y=y,J=j]
 &=T\xi_j+TC_jS^\top(R_j^Y)^{-1}(y-S\xi_j),\\
 \Cov(V\mid Y=y,J=j)
 &=TC_jT^\top-TC_jS^\top(R_j^Y)^{-1}SC_jT^\top.
 \label{eq:correction}
\end{align}
Thus a predictive value $T\xi_j$ without the correction is exact for every observation only when $TC_jS^\top=0$. Euclidean row orthogonality is insufficient. For $S=(1,0)$, $T=(0,1)$ and covariance with off-diagonal entry $\rho$, the correction is $\rho(y-\xi_{j,1})$ when the first marginal variance is one and observation noise is absent.

The implemented predictive readout uses the separate conditional-content interpretation. It does not claim to condition a shared latent Gaussian draw and does not calculate \eqref{eq:correction}. Independent learned maps alone cannot make these two probability laws equivalent.

\section{Separate predictive-content readout with general output mixing}
\subsection{The returned mean}
Define each source value and each pooled head content by
\begin{equation}
 v_j^h=T_h\xi_j^K,\qquad c_i^h=\sum_j\beta_{ij}^h v_j^h.
\end{equation}
The canonical output and its query-frame realization are
\begin{equation}
 \overline m_i=\sum_{h=1}^H O_h c_i^h,\qquad
 \boxed{m_i=A_i\sum_{h=1}^H O_h\sum_j\beta_{ij}^h T_hB_j^{-1}\mu_j.}
 \label{eq:main}
\end{equation}
No pairwise $K$-dimensional transported value array is required: form the compact source content, pool it with attention, apply the output maps, and finally apply the receiver frame. Independent query/key frames enter the same construction coherently through $A_i$ and $B_j$.

Let $c_i$ concatenate the $c_i^h$ and let $W_O=[O_1\ \cdots\ O_H]$. Then $\overline m_i=W_Oc_i$. This is exactly the column-vector form of concatenation followed by an arbitrary linear output projection. The old block-output family imposes $O_h=E_h\overline O_h$ with $\overline O_h\in\R^{d\times s_h}$. General output mixing removes that row-support restriction. It therefore adds a different freedom from whole-belief input access.

A dense output map also differs from adding off-diagonal gauge generators. It changes how readout content is written into a canonical representation, while leaving the existing compact frame algebra intact. It does not require merging the gauge blocks into a single attention head. Likewise, a constrained post-attention HeadMixer is not the definition of a general $W_O$. When general O is enabled, the existing HeadMixer may still be used as an additional transformation, but it is not necessary to provide unrestricted output mixing.

\subsection{Predictive covariance and what the implementation does not return}
If head $h$ has source-conditioned content covariance $R_j^h$, its mixture covariance is
\begin{equation}
 V_i^h=\sum_j\beta_{ij}^h\left[R_j^h+
 (v_j^h-c_i^h)(v_j^h-c_i^h)^\top\right].
 \label{eq:mixturecov}
\end{equation}
This is the within-component covariance plus the between-component covariance. If head-content draws are independent conditional on the routing context, the output covariance is
\begin{equation}
 \Sigma_i^{\rm pred}=A_i\left(\sum_h O_hV_i^hO_h^\top\right)A_i^\top.
\end{equation}
Without that additional independence assumption, cross terms $O_h\Cov(Z_i^h,Z_i^l)O_l^\top$ are also needed. Overlapping measurements of one latent belief do not justify dropping them.

The present software computes only \eqref{eq:main}. At the readout seam it carries the receiver covariance forward unchanged to satisfy the existing belief-state interface. That stored covariance is bookkeeping, not \eqref{eq:mixturecov} and not a propagated T/O covariance. Later configured normalization, HeadMixer, or block-MLP transformations retain their existing behavior. The returned data structure must not be interpreted as a newly derived joint predictive Gaussian solely because its mean came from \eqref{eq:main}.

\subsection{Predictive averaging versus a belief-coordinate update}
A Gaussian mean update against heteroscedastic content densities has precision weighting. For fixed weights, means $v_j$, and precisions $\Lambda_j\succ0$,
\begin{equation}
 E(m)=\frac12\sum_j\beta_j(m-v_j)^\top\Lambda_j(m-v_j)
\end{equation}
has minimizer
\begin{equation}
 m_*=\left(\sum_j\beta_j\Lambda_j\right)^{-1}
       \sum_j\beta_j\Lambda_jv_j.
\end{equation}
This is generally different from $\sum_j\beta_jv_j$. A finite Fisher-preconditioned step toward this minimizer also depends on the starting mean, step size, preconditioner, and any trust bound or additional self term. Equality to direct predictive averaging requires specific conditions, not merely the existence of learned value maps.

For a separate content state with common covariance $R$, the metric on the mean is $R^{-1}$ and the natural gradient of this fixed-context quadratic is $m-\sum_j\beta_jv_j$. A unit step reaches the predictive mean. This is a conditional optimization interpretation of a separate state. It does not prove that the comparison query itself follows that flow.

To see the obstruction, take query $q\in\R^2$, keys $\pm(1,0)^\top$, and values $\pm(0,1)^\top$. The attention increment is $f(q)=(0,\tanh q_1)^\top$. At zero its Jacobian is $N=\left(\begin{smallmatrix}0&0\\1&0\end{smallmatrix}\right)$. If $f=-G^{-1}\nabla E$ for a smooth positive-definite metric and a twice differentiable energy, then $f(0)=0$ implies $\nabla^2E(0)=-G(0)N$. Writing $G(0)=\left(\begin{smallmatrix}a&b\\b&c\end{smallmatrix}\right)$, symmetry of the Hessian forces $c=0$, contradicting positive definiteness. Arbitrary independent-value attention increments therefore need not be natural-gradient flows on the same query coordinates. This local example concerns increments and does not forbid the separate predictive map.

\section{Gauge covariance and role-frame freedom}
Under a passive local coordinate change $g_i\in GL(K)$, transform all local quantities together:
\begin{equation}
 \mu_i'=g_i\mu_i,\quad\Sigma_i'=g_i\Sigma_i g_i^\top,\quad
 A_i'=g_iA_i,\quad B_i'=g_iB_i.
\end{equation}
Equation \eqref{eq:canonical} is unchanged, so the canonical comparison marginals, routing weights, and content means are unchanged. Applying \eqref{eq:main} yields $m_i'=g_im_i$. The T/O matrices need not commute with every $g_i$, because they act between explicitly chosen canonical readout spaces, rather than acting as fixed local-coordinate endomorphisms.

This algebraic covariance statement assumes that the covariance family and frame representation support the stated transformation. Arbitrary $GL(K)$ transformations do not preserve a diagonal local covariance family; arbitrary dense transformations do not preserve a block-diagonal frame representation. Thus the implementation's diagonal belief approximation does not establish full-model $GL(K)$ equivariance. It computes the exact reduced marginal of the approximate Gaussian it actually stores. The strict positive-definiteness requirement, finite-precision policy, and any downstream normalization are additional parts of the executable model.

With separate role frames, the pair map is $A_iB_j^{-1}$. It need not satisfy a flat same-frame cocycle: multiplying $A_iB_j^{-1}$ by $A_jB_k^{-1}$ leaves the middle factor $B_j^{-1}A_j$. Tying $A_j=B_j$ removes it. The predictive construction requires coherent role use, not a false assertion that separate-role transport is flat. Additional role frames enlarge the architecture beyond the standard four learned projection matrices.

\section{Capacity, initialization, and optimization}
For equal ranks $r,s$ and $H$ heads, independent learned comparison and content parameters have the following counts; frame, embedding, normalization, decoder, and other block parameters are excluded:
\begin{center}
\begin{tabular}{lll}
\toprule
Input and output support & Q/K parameters & T/O parameters\\
\midrule
Head input, head output & $2Hdr$ & $2Hds$\\
Full input, head output & $2HKr$ & $Hs(K+d)$\\
Full input, full output & $2HKr$ & $2HKs$\\
\bottomrule
\end{tabular}
\end{center}
For $r=s=d=K/H$, the last row totals $4K^2$, the usual four-map parameter count. At $K=20,H=2,r=s=10$, the three totals are respectively $800$, $1400$, and $1600$. Shared or fixed comparison maps and identity content have different counts. A capacity increase is a possible experimental confound, not evidence that the architectural change will help.

The implementation initializes full-output O by embedding each existing local O into its old output block and setting the other output rows to zero. Those rows remain trainable. With $s=d$, T/O therefore begin at local identity placement; smaller $s$ begins at a local projector. This initial support does not constrain the learned support. Independent Q/K matrices begin with orthonormal rows but train as free matrices, so their singular values and reduced-covariance condition numbers remain relevant diagnostics.

The predictive path keeps outer gradients through source means, comparison covariances, Q/K maps, T/O maps, and realized frames. In particular, reading a source value does not detach the source belief. That is a deliberate difference from the older filtering path's detached source context. The common phrase ``fixed context'' in the conditional model refers to how the forward probability law is defined; it does not instruct the outer optimizer to sever that law's dependence on learned context parameters.

\section{Executable selectors and the closest Transformer-like setting}
\subsection{The implemented branch}
The two selectors are \texttt{relational\_readout\_mode} and \texttt{relational\_output\_scope}. The first accepts \texttt{conditional\_update} or \texttt{predictive\_content}; the second accepts \texttt{head} or \texttt{full}. Existing defaults remain the conditional update with head output. Full output requires predictive content and learned S/T/O operators. With the relational router, predictive content requires an integer learned content rank. The legacy \texttt{projected\_filtering} mode name selects the relational comparison infrastructure even when the readout selector replaces its filtering update. With relational mode off and all legacy QKV map flags disabled, the predictive selector instead performs parameter-free averaging using the unprojected router.

In the predictive branch, the E-step seam evaluates \eqref{eq:main} directly and replaces the mean for that pass. It does not use the old mean learning rate, mean Fisher step, mean trust bound, self-attraction force, or value-energy clipping. It adds no attention residual connection. The same source mask, KL comparison clamp, and configured temperature remain in force. Covariance is frozen at this seam, and the existing downstream block pipeline continues. Conditional-cost fields and legacy global-free-energy fields are explicitly unavailable, while routing, displacement, operator-rank, parameter-count, and comparison-conditioning diagnostics remain meaningful.

The first implementation keeps the established conditional envelope: one layer, one E-step, diagonal Gaussian local beliefs, KL comparison, frozen covariance and E-step frames, compact flat block-$GL$ frame construction, and the supported linear decoding path. Full source Gaussian families, arbitrary depth, and a replacement dot-product score are not silently enabled by these selectors. Configuration schema 19 stores the new fields; older checkpoints migrate to the prior conditional/head behavior.

\subsection{Predictive content with unprojected routing}
The following unprojected alternative isolates the readout choice from the introduction of learned S/T/O maps. It is useful when comparing update laws within the existing Gaussian-belief architecture.
\begin{verbatim}
relational_mode                  = "off"
relational_readout_mode          = "predictive_content"
relational_input_scope           = "full"  # or "head"
relational_output_scope          = "head"
use_w_q                         = False
use_w_k                         = False
use_w_v                         = False
relational_score_rank_per_head   = None
relational_content_rank_per_head = None
\end{verbatim}
For transported means $u_{ij}=\Omega_{ij}\mu_j$, this branch returns
\begin{equation}
 (m_i)^{(h)}=\sum_j\beta_{ij}^h(u_{ij})^{(h)}.
\end{equation}
Here $\Omega_{ij}=A_iB_j^{-1}$ when separate query and key frames are selected, and the superscript selects the runtime output block. Head-local scope uses the existing blockwise unprojected KL score. Full scope uses one whole-belief KL energy shared across heads before applying their priors and temperatures. Scoring retains the existing diagonal readout of the transported covariance; it does not use an S-projected Gaussian marginal. Both scopes keep source and frame gradients live, bypass the old inference mean step, and retain receiver covariance as bookkeeping. This branch adds zero readout parameters. Full learned output mixing still requires the S/T/O router, and predictive use of the legacy learned QKV mean maps remains outside the implemented scope.

\subsection{Closest supported Transformer map and readout structure}
At the current example width $K=20,H=2$, the following settings choose the closest supported whole-input and general-output construction:
\begin{verbatim}
relational_mode                  = "projected_filtering"
relational_input_scope           = "full"
relational_qk_mode               = "independent"
relational_score_rank_per_head   = 10
relational_content_rank_per_head = 10
relational_readout_mode          = "predictive_content"
relational_output_scope          = "full"
qk_frame_mode                    = "tied"
use_head_mixer                   = False
use_w_q                         = False
use_w_k                         = False
use_w_v                         = False
\end{verbatim}
The last three switches disable the separate legacy mean-only QKV adapter. They do not disable the learned relational Q/K/T/O maps. The tied-frame choice avoids adding distinct query and key gauge fields beyond the four-map architecture. Turning off HeadMixer removes an additional constrained mixing stage now that O itself has general output support. Other required envelope settings must still be retained; this block is an override to a valid relational configuration, not a complete training configuration.

The map dimensions and direct weighted-value readout agree with standard multi-head attention. The score still differs. Even with identical isotropic comparison covariances, Gaussian KL produces squared distance and hence a source key-norm penalty when expanded. With source-dependent covariances it also has source-dependent precision, trace, and log-determinant terms. The software temperature remains $\kappa_\beta\sqrt d$, rather than automatically changing to $\sqrt r$ when rank changes. Setting $r=d$ aligns those dimensions, but does not convert KL to dot product. Exact vanilla attention additionally requires the appropriate score, treatment of frames, residual connections, normalization, feed-forward layers, masking, and training objective. No available selector combination should be described as an exact complete Transformer on the basis of this change alone.

\section{Tests and experiments that answer different questions}
Small numerical tests can establish that the executable contraction matches an explicit dense reference, that gradients reach all four map families, that off-block output rows are trainable, and that switching the old E-step mean learning rate does not change the predictive branch. Checkpoint, evaluation, and figure tests establish that the new semantics survive the surrounding run lifecycle. They do not establish an accuracy advantage.

A controlled experiment should distinguish the input-support change, the readout-law change, and the output-support change. Holding full input, independent Q/K, ranks, frames, data order, and training budget fixed, compare the conditional/head-output path with predictive/head output, then predictive/full output. The first comparison changes both the mean-update law and the source-gradient convention in this implementation, so it is not a pure test of precision weighting alone. The second changes output support and parameter count. Comparing head input with full input adds another capacity difference. A separate parameter-matched comparison can address that difference, but necessarily changes ranks or another resource and must state which.

Use repeated matched seeds, record completed optimization steps and tokens as well as wall time, and report uncertainty across runs. Freeze or record every relevant configuration and checkpoint revision. A one-seed screening result is descriptive. No training sweep or production GPU experiment was run for this implementation, and no mechanism has been established as the cause of any earlier perplexity plateau.

\section{Remaining modeling choices}
The separate predictive variable resolves the need to interpret arbitrary weighted values as a force on the comparison query. It leaves several modeling choices explicit: whether predictive uncertainty should be represented, whether head-content variables are conditionally independent, whether source selection should be revised after observing a target, and whether a future multilayer model should maintain a comparison belief alongside predictive content rather than reuse one state container. Those choices require their own model and interface contracts.

The construction already supplies an executable architectural test: learned whole-belief Q/K and T maps determine comparison and content, while general O maps distribute pooled head content over the output space. Its exact mathematical claims concern the displayed conditional laws, moments, and coordinate transformations. Its current software claim concerns a mean-only predictive readout within a bounded existing runtime. Predictive performance remains an empirical question.

\begin{thebibliography}{9}
\bibitem{vaswani} A. Vaswani et al., ``Attention Is All You Need,'' Advances in Neural Information Processing Systems 30 (2017), Sections 3.2.1--3.2.2. \href{https://arxiv.org/abs/1706.03762}{arXiv:1706.03762}.
\end{thebibliography}
\end{document}
